\section{ET-MCP Architecture}
\label{sec:architecture}

ET-MCP is the substrate that makes the measurement protocol of
\S\ref{sec:evaluation} possible; it is not the contribution of
this paper. This section covers only what the measurement story
needs: the minimal tool surface, the failure-only writer policy,
and the pull-vs-intercept reader-side distinction that
\S\ref{sec:evaluation} compares.

\subsection{Tool surface and lifecycle}
\label{sec:arch:tools}
\label{sec:arch:lifecycle-write-query}

The substrate is structured around one trace namespace per
\texttt{task\_id}, served over the stateless MCP 2026-07-28
transport~\cite{mcpspec2026}. Agents see three operations:
\texttt{trace.write} (append a typed event),
\texttt{trace.query} (TF-IDF ranking over the typed payload,
returning top-$k$ peer events with one-line summaries), and
\texttt{trace.cas} (compare-and-swap for the rare double-commit
race). The task lifecycle (\texttt{init}, \texttt{complete}, TTL
sweep) is owned by the orchestrator through a separate channel
that participating agents cannot reach, so no agent can tear down
a peer's namespace. The TF-IDF ranker is a single replaceable
component.

\subsection{Writer policy}
\label{sec:arch:taxonomy}

Events carry an envelope (\texttt{event\_id}, \texttt{task\_id},
\texttt{event\_type}, \texttt{agent\_id}, \texttt{timestamp},
\texttt{version}, \texttt{payload}). Five typed event categories
cover the negative-outcome cases mapped to the MAST failure
taxonomy~\cite{mast2025}: failed paths, constraint violations,
abandoned approaches, irreversible intermediate decisions, and
tool errors. The default \texttt{failure\_only} client-side
selection policy writes only on these negative outcomes. The
$\tau^2$ retail evaluation specializes this to a single
\textsc{failed\_trial\_action} type keyed on the trial's terminal
reward.

\subsection{Two reader-side architectures: pull and intercept}
\label{sec:arch:pull-vs-intercept}

\begin{figure*}[t]
\centering
% Side-by-side architecture diagram comparing pull and intercept.
% Two-column figure (figure* in 04-architecture.tex).
% Semantic color encoding consistent across panels:
%   blue   = store read (curated summary / consult)
%   gray   = ordinary agent <-> tool I/O
%   red    = framework-injected / peer-warning payload (intercept only)
% Step labels are numbered per panel (P1..P3 / I1..I5); the caption
% decodes every step in prose so the figure itself stays uncluttered.
% Requires: \usetikzlibrary{positioning,arrows.meta,fit,backgrounds,shapes.geometric,calc}
\begin{tikzpicture}[
  font=\small,
  agent/.style={draw, rectangle, rounded corners=2pt,
                minimum width=1.7cm, minimum height=0.7cm,
                inner sep=2pt, fill=blue!10, align=center},
  tool/.style={draw, rectangle, rounded corners=2pt,
               minimum width=1.7cm, minimum height=0.7cm,
               inner sep=2pt, fill=green!10, align=center},
  framework/.style={draw, rectangle, rounded corners=2pt, dashed, thick,
                    minimum width=2.4cm, minimum height=0.7cm,
                    inner sep=2pt, fill=orange!12, align=center},
  store/.style={draw, cylinder, shape border rotate=90, aspect=0.25,
                minimum width=2.0cm, minimum height=0.8cm,
                inner sep=2pt, fill=gray!12, align=center, font=\footnotesize},
  flow/.style={-{Stealth[length=2mm,width=1.8mm]}, thick, draw=black!65, line width=0.6pt},
  flowread/.style={-{Stealth[length=2mm,width=1.8mm]}, thick, draw=blue!75!black, line width=0.75pt},
  flowwarn/.style={-{Stealth[length=2mm,width=1.8mm]}, thick, draw=red!75!black, line width=0.85pt, dashed},
  stepnum/.style={font=\bfseries\scriptsize, circle, draw=black!70, inner sep=0.4pt,
               minimum size=3.6mm, fill=white},
  stepwarn/.style={font=\bfseries\scriptsize, circle, draw=red!70!black, inner sep=0.4pt,
                   minimum size=3.6mm, fill=red!8},
  stepread/.style={font=\bfseries\scriptsize, circle, draw=blue!70!black, inner sep=0.4pt,
                   minimum size=3.6mm, fill=blue!6},
  panel/.style={draw=black!60, thick, rounded corners=4pt, inner sep=10pt, fill=white}
]

% =========================================================
% LEFT PANEL: PULL  (agent-as-reader)
% Layout: store top, agent + tool side-by-side below.
% =========================================================
\begin{scope}[local bounding box=pullbox]
  \node[store]                    (s1) at (0, 2.2) {ET-MCP store};
  \node[agent, below=14mm of s1]  (a1)              {Agent};
  \node[tool, right=22mm of a1]   (t1)              {Tool};

  % (P1) store -> agent: curated summary at trial start
  \draw[flowread] (s1.south) -- (a1.north)
    node[stepread, pos=0.45] {P1};

  % (P2) agent -> tool: tool call (upper lane)
  \draw[flow]
    ([yshift=3mm]a1.east) -- ([yshift=3mm]t1.west)
    node[stepnum, pos=0.5] {P2};

  % (P3) tool -> agent: raw response (lower lane)
  \draw[flow]
    ([yshift=-3mm]t1.west) -- ([yshift=-3mm]a1.east)
    node[stepnum, pos=0.5] {P3};
\end{scope}
\begin{scope}[on background layer]
\node[panel, fit=(pullbox),
      label={[font=\bfseries]above:Pull \emph{(agent-as-reader)}}] (pullpanel) {};
\end{scope}

% =========================================================
% RIGHT PANEL: INTERCEPT  (framework-as-reader)
% Layout: store top, framework below store, agent + tool side-by-side below framework.
% Curved arrows so request (top) and response (bottom) lanes are clearly separable.
% =========================================================
\begin{scope}[xshift=10cm, local bounding box=icbox]
  \node[store]                       (s2) at (0, 2.5) {ET-MCP store};
  \node[framework, below=8mm of s2]  (fw)             {Framework};
  \node[agent, below=22mm of fw, xshift=-22mm]   (a2) {Agent};
  \node[tool,  below=22mm of fw, xshift=22mm]    (t2) {Tool};

  % I4: framework consults store (vertical, straight)
  \draw[flowread] (fw.north) -- (s2.south)
    node[stepread, pos=0.5] {I4};

  % I1: agent -> framework (upper curve, bending right/up)
  \draw[flow, bend left=20]
    (a2.north) to
    node[stepnum, pos=0.55] {I1}
    ([xshift=-3mm]fw.south);
  % I5: framework -> agent (lower curve, bending right/down, red dashed)
  \draw[flowwarn, bend left=20]
    ([xshift=-6mm]fw.south) to
    node[stepwarn, pos=0.45] {I5}
    ([xshift=3mm]a2.north);

  % I2: framework -> tool (upper curve, bending left/down)
  \draw[flow, bend left=20]
    ([xshift=3mm]fw.south) to
    node[stepnum, pos=0.45] {I2}
    (t2.north);
  % I3: tool -> framework (lower curve, bending left/up)
  \draw[flow, bend left=20]
    ([xshift=-3mm]t2.north) to
    node[stepnum, pos=0.55] {I3}
    ([xshift=6mm]fw.south);
\end{scope}
\begin{scope}[on background layer]
\node[panel, fit=(icbox),
      label={[font=\bfseries]above:Intercept \emph{(framework-as-reader)}}] (icpanel) {};
\end{scope}

% Compact semantic legend below both panels (single line).
\node[font=\scriptsize, below=2.5mm of pullpanel.south west, anchor=north west,
      xshift=2mm] {%
  \tikz[baseline=-0.5ex]{\draw[flowread] (0,0)--(0.55,0);}\;store read\quad
  \tikz[baseline=-0.5ex]{\draw[flow] (0,0)--(0.55,0);}\;agent$\leftrightarrow$tool\quad
  \tikz[baseline=-0.5ex]{\draw[flowwarn] (0,0)--(0.55,0);}\;injected \texttt{[PEER-WARNING]}};

\end{tikzpicture}

\caption{Reader-side architecture contrast. The two designs share
the event store, writer policy, and MCP transport; they differ
only in who reads and when the payload is presented.
\textbf{Pull} (left): P1 store $\to$ agent (curated summary at
trial start); P2 agent $\to$ tool (call); P3 tool $\to$ agent
(raw response).
\textbf{Intercept} (right): I1 agent $\to$ framework (call); I2
framework $\to$ tool (forward); I3 tool $\to$ framework (raw);
I4 framework $\to$ store (consult); I5 framework $\to$ agent
(response, with \texttt{[PEER-WARNING]} prepended on argument
match).  Pull presents the curated payload once; intercept
injects raw warnings per matching call.}
\label{fig:pull-vs-intercept}
\Description{Side-by-side architecture diagram. Left panel shows
the pull architecture: the ET-MCP store sends a curated summary to
the agent at trial start, then the agent calls the tool and the
tool returns a raw response. Right panel shows the intercept
architecture: the agent calls the tool, the call passes through
the framework which consults the ET-MCP store, the tool returns
its raw response to the framework, and the framework forwards a
modified response with a [PEER-WARNING] block prepended back to
the agent.}
\end{figure*}

Figure~\ref{fig:pull-vs-intercept} contrasts the two designs.
\textbf{Pull} places the agent as the \emph{active reader}: it
calls \texttt{trace.query} when it estimates peer events might
inform the next decision, paying decision overhead, round-trip
latency, and an enlarged tool surface for that control.

\textbf{Intercept} relocates the reader into the framework. The
orchestrator intercepts every tool call, and if the call's
(name, arguments) appears in the task-scoped store as a peer
\textsc{failed\_path} or \textsc{failed\_trial\_action} event,
it prepends a \texttt{[PEER-WARNING]} block to the response.
Intercept removes the agent-side decision, the LLM round-trip,
and any agent code changes; the tool surface and prompt are
unchanged. The hook surface intercept rides on is already a
production primitive in Strands~\cite{strandshooks2025},
AgentCore Policy~\cite{agentcorepolicy2025}, LangChain
middleware~\cite{langchainmiddleware2026}, and Microsoft Agent
Framework~\cite{msagentframework2026}, so intercept is an
additive extension to existing runtimes, not a new framework.

Both architectures share the same event taxonomy, failure-only
writer policy, task-scoped lifecycle, and MCP transport, and are
configuration-equivalent at trial 0: pull's augmenter is a no-op
when the store is empty, and the intercept hook is pass-through
on no argument match. The full audit of this equivalence is in
\S\ref{sec:eval:request-equivalence}. The next
section instantiates both in the reference implementation that
produces the measurements of \S\ref{sec:evaluation}.
